\documentclass[11pt,a4paper]{article}
\usepackage[utf8]{inputenc}
\usepackage[T1]{fontenc}
\usepackage{amsmath,amssymb,amsthm}
\usepackage{mathtools}
\usepackage{hyperref}
\usepackage{geometry}
\usepackage{xcolor}
\usepackage{booktabs}
\usepackage{enumitem}
\geometry{margin=1in}

\hypersetup{colorlinks=true, linkcolor=blue!50!black, urlcolor=blue!50!black,
  citecolor=blue!50!black}

%% =====================================================================
%% Standalone applied-ML note (2026-06-23).  Reports the learnable per-head
%% positional bias selected by elimination across three GPU bake-offs, its
%% length-extrapolation behaviour, and a careful (illustrative, NOT measured)
%% discussion of what stable train-short/serve-long extrapolation could mean
%% for inference-time energy and CO2.  Modest, falsification-first, forward
%% looking.  Companion to the Mirror-Map-Sieve preprint (v7).
%% =====================================================================

\newtheorem{observation}{Observation}

\title{Adaptability Beats Added Structure:\\
  A Learnable Per-Head Positional Bias, Selected by Elimination,\\
  and What Its Length Extrapolation Could Mean for the Energy\\
  Footprint of Long-Context Inference\\
  \large (applied note --- exploratory, seeking review)}
\author{Xavier Callens\\
  SocrateAI Laboratory\\
  \url{https://github.com/xaviercallens/Mirror-Map-Sieve}}
\date{\today}

\begin{document}
\maketitle

\begin{abstract}
We report a small, deliberately falsifiable applied-ML result and discuss its
forward-looking implications without overclaiming. Starting from a geometric
prior (a positional-attention bias derived from a Calabi--Yau--type period), we
ran a sequence of from-scratch training comparisons on WikiText-2 and were led,
by a chain of negative and null results, to a simple conclusion: among a family
of increasingly expressive per-head positional biases, the one that wins at GPU
scale is the \emph{plainest}---a single \emph{learnable} linear slope per
attention head (``learnable-ALiBi''), with no log-curvature term, no
number-theoretic sequence, no per-head content scale, and no per-head softmax
temperature. Three independent head-to-head studies (an attribution experiment,
a six-family bake-off, and a content/position-balance probe) converge on this
scheme; every added degree of freedom we tested either ties it within noise or
\emph{regresses} once the model is properly trained and regularized. The scheme
PASSes a pre-committed quality gate ($+8.1\%$ perplexity over the best baseline
on a regularized $8000$-step run) and---the property we find most
interesting---\emph{extrapolates flat} in sequence length, degrading gracefully
out to $8\times$ the training context where learned absolute positions collapse.
We are explicit that this is a known-good idea (ALiBi with learnable slopes),
modestly improved and carefully validated, not a new mechanism. We then discuss,
as \emph{illustrative order-of-magnitude reasoning and not as a measured
result}, why robust train-short / serve-long extrapolation is precisely the kind
of property that bears on the energy and $\mathrm{CO_2}$ cost of LLM serving at
scale, and we lay out a concrete, fair-comparison roadmap (vs.\ NoPE, RoPE +
position-interpolation, and learnable-ALiBi on a genuine long-range benchmark)
that would be needed to substantiate any such claim. Code, raw runs, and an open
benchmark dataset accompany this note.
\end{abstract}

\tableofcontents

\section{Scope, and what we do \emph{not} claim}
\label{sec:scope}

This is an applied side-note to a mathematics preprint~\cite{mms}, not a core
result. To keep the reader's expectations calibrated, we state our limits up
front.

\begin{itemize}[leftmargin=1.4em,itemsep=2pt]
\item The mechanism we end on---\textbf{a learnable per-head linear positional
  bias}---is \emph{not new}. It is ALiBi~\cite{alibi} with the per-head slopes
  made trainable rather than fixed. Our contribution is empirical and
  methodological: a controlled elimination showing that, on this benchmark,
  nothing more elaborate survives proper training, together with an honest
  measurement of its length-extrapolation behaviour.
\item All quality numbers are on \textbf{WikiText-2}~\cite{wikitext}, a single
  small corpus, with models trained \emph{from scratch} at the $25$M-parameter
  scale. They are suggestive, not conclusive; they should be reproduced at
  larger scale and on more corpora before being relied upon.
\item We claim \textbf{no memory or FLOP advantage} over a modern kernel. A
  learnable per-head scalar bias is $\mathcal{O}(1)$ state and $\mathcal{O}(L)$
  to apply---identical to ALiBi as implemented inside
  FlashAttention~\cite{flash}. Earlier framings of this project that suggested a
  memory ``breakthrough'' were against a materialized-table strawman and are
  retracted.
\item The energy / $\mathrm{CO_2}$ discussion in
  Section~\ref{sec:energy} is \textbf{explicitly illustrative}. We did not
  measure datacenter power. We reason about where a robust-extrapolation
  positional scheme \emph{could} reduce work, and we are careful to separate that
  reasoning from evidence.
\end{itemize}

\noindent The reason we publish at all is the discipline: a bold prior was
falsified, a false positive was caught, a modest true result was confirmed, and a
controlled experiment then corrected our own story about \emph{why} it works. We
think that arc is worth recording even though the surviving mechanism is humble.

\section{Background: positional bias as an additive score}
\label{sec:bg}

In a decoder attention layer with head $h$, query $i$ and key $j\le i$, the
pre-softmax score is
\begin{equation}
  s_h(i,j) \;=\; \frac{q_i^\top k_j}{\sqrt{D}} \;+\; \beta_h(i,j),
  \qquad j\le i,
  \label{eq:score}
\end{equation}
where $\beta_h$ is an additive positional bias and $D$ the head dimension.
Several widely used schemes are special cases of~\eqref{eq:score}:
\begin{align}
  \text{learned-absolute:}\quad & \beta_h(i,j)=0,\ \text{positions added to embeddings};\\
  \text{ALiBi~\cite{alibi}:}\quad & \beta_h(i,j) = -a_h\,(i-j),\ a_h\ \text{a \emph{fixed} per-head slope};\\
  \text{NoPE~\cite{nope}:}\quad & \beta_h(i,j)=0,\ \text{no positional signal at all}.
\end{align}
ALiBi's fixed geometric slope ladder $a_h = 2^{-8h/H}$ extrapolates well but is
not adapted to the data; learned-absolute fits well in-distribution but cannot
represent positions beyond the training length, so it collapses on
extrapolation. The question that organizes this note is whether a \emph{more
expressive but still cheap} $\beta_h$ can keep ALiBi's extrapolation while
improving in-distribution quality---and, if so, \emph{how much} expressiveness
actually helps.

\section{Method and protocol}
\label{sec:method}

\paragraph{The candidate family.} We consider per-head biases of increasing
expressiveness, all $\mathcal{O}(1)$ in stored parameters per head:
\begin{align}
  \text{(slope only)}\quad &\beta_h(d) = -\,a_h\,d, & a_h&\ \text{learnable};\label{eq:alibilearn}\\
  \text{(slope $+$ log)}\quad &\beta_h(d) = -\,a_h\,d + b_h\log d, & a_h,b_h&\ \text{learnable};\label{eq:nested}\\
  \text{(content balance)}\quad & s_h(i,j) = c_h\frac{q_i^\top k_j}{\sqrt D} - a_h\,d, & a_h,c_h&\ \text{learnable};\label{eq:cb}\\
  \text{(}+\text{temperature)}\quad & s_h(i,j) = \tau_h^{-1}\Big(c_h\frac{q_i^\top k_j}{\sqrt D} - a_h\,d\Big), & a_h,c_h,\tau_h&\ \text{learnable},\label{eq:cbtemp}
\end{align}
with $d=i-j$. Equation~\eqref{eq:alibilearn} is \emph{learnable-ALiBi}, our
eventual selection. Equation~\eqref{eq:nested} adds a learnable log-curvature
(this is also where a fixed Calabi--Yau curve, $\beta_h(d)=-\gamma_h\log S_{20}(d)$,
sits as a constrained special case). Equation~\eqref{eq:cb} introduces a
per-head \emph{content scale} $c_h$ that lets a head trade off content matching
against positional decay---making the ``local head vs.\ global head'' split
explicit rather than emergent. Equation~\eqref{eq:cbtemp} adds a per-head softmax
temperature.

\paragraph{Pre-committed gate.} Following the ALiBi evaluation
methodology---train one model \emph{from scratch} per positional scheme at
matched compute---we fixed, in advance, a quality gate: \textsc{pass} within
$+1\%$ of the best baseline validation perplexity, \textsc{kill} at $>5\%$
regression. For each more-expressive scheme we also pre-committed a
\emph{specific} null: e.g.\ ``content balance must improve length extrapolation
over learnable-ALiBi, or it is killed.''

\paragraph{Training.} Decoder-only Transformers, $d_{\text{model}}=512$, $8$
layers, $8$ heads ($\approx 25$M parameters), trained on real WikiText-2
(\texttt{Salesforce/wikitext}, \texttt{wikitext-2-raw-v1}) on a single NVIDIA L4,
PyTorch~2.9.1\,/\,CUDA~12.9. The decisive runs use $4000$ steps, context $512$,
batch $24$, with $\ell_2$ regularization on the learnable slopes ($\lambda=10^{-3}$)
and validation early-stopping---both of which proved \emph{essential} (see
Section~\ref{sec:lessons}). Length extrapolation is measured as validation
perplexity at $1\times,2\times,4\times,8\times$ the training context, with no
retraining.

\section{Results}
\label{sec:results}

\subsection{Attribution: the gain is learnability, not the geometric shape}
A controlled nested experiment~\eqref{eq:nested} separates the linear slope from
the curvature term (Table~\ref{tab:attr}). A plain learnable slope with
\emph{zero} curvature already beats a fixed geometric (Calabi--Yau) bias; and
when the curvature $b_h$ is left free starting from the geometric value, gradient
descent drives it \emph{away} from that value (toward and past zero). The
geometry was a prior, and not a well-pointed one; what helps is letting the data
choose the slope.

\begin{table}[t]
\centering
\caption{Attribution (local controlled run). Perplexity vs.\ fixed ALiBi; lower
ppl is better. A learnable slope with no curvature (\emph{alibi\_learn}) already
beats the fixed geometric shape (\emph{holo\_fixed}).}
\label{tab:attr}
\begin{tabular}{lcc}
\toprule
scheme & ppl @512 & vs.\ fixed ALiBi \\
\midrule
nested (slope $+$ free log, linear init) & 8.548 & $+6.9\%$ \\
nested (slope $+$ free log, geometric init) & 8.680 & $+5.5\%$ \\
\textbf{learnable-ALiBi (slope only)} & 8.835 & $+3.8\%$ \\
fixed geometric (Calabi--Yau) shape & 8.998 & $+2.0\%$ \\
ALiBi (fixed slopes) & 9.181 & --- \\
\bottomrule
\end{tabular}
\end{table}

\subsection{Scale check: the win is real and PASSes the gate}
On a regularized $8000$-step run, the learnable scheme reaches $3.324$ validation
perplexity versus $3.617$ for the best baseline (learned-absolute)---a
\textbf{$+8.1\%$ improvement that PASSes} the pre-committed gate, with the
overfitting that plagued an earlier under-regularized run (Section~\ref{sec:lessons})
removed.

\subsection{Elimination: nothing more expressive survives}
The decisive comparison~(Table~\ref{tab:hetero}) trains the
schemes~\eqref{eq:alibilearn}--\eqref{eq:cbtemp} from scratch and measures length
extrapolation. The explicit content-balance knob $c_h$ \eqref{eq:cb} only
\emph{ties} learnable-ALiBi ($+0.35\%$ at $4\times$, $+0.17\%$ at $8\times$---within
noise), and adding the per-head temperature $\tau_h$ \eqref{eq:cbtemp}
\emph{regresses} at every length. The content-only control (NoPE) collapses and,
tellingly, gets \emph{worse} with length, confirming the positional term is doing
real work.

\begin{table}[t]
\centering
\caption{Length extrapolation (NVIDIA L4, from scratch, ctx $512$, $4000$ steps).
Validation perplexity at $1$--$8\times$ the training context; lower is better.
The added content scale ties learnable-ALiBi; the temperature knob regresses.}
\label{tab:hetero}
\begin{tabular}{lcccc}
\toprule
scheme & $1\times$ (512) & $2\times$ (1024) & $4\times$ (2048) & $8\times$ (4096) \\
\midrule
ALiBi (fixed) & 3.705 & 3.623 & 3.577 & 3.513 \\
\textbf{learnable-ALiBi} & \textbf{3.651} & \textbf{3.559} & 3.500 & 3.429 \\
$+$ content scale $c_h$ & 3.642 & 3.548 & \textbf{3.488} & \textbf{3.428} \\
$+$ content $+$ temperature & 3.732 & 3.659 & 3.599 & 3.548 \\
NoPE (content only) & 10.290 & 10.845 & 11.571 & 12.311 \\
\bottomrule
\end{tabular}
\end{table}

\begin{observation}
Across all three studies, expressiveness beyond a single learnable per-head slope
does not survive proper training at this scale: it either ties within noise or
regresses. We summarize this as \emph{adaptability beats added structure---up to
one learnable slope per head, and no further}.
\end{observation}

\subsection{The property worth keeping: graceful extrapolation}
The learnable scheme not only edges the baselines in-distribution; its perplexity
is \emph{flat-to-improving} from $1\times$ to $8\times$ the training context
(Table~\ref{tab:hetero}), whereas learned-absolute positions degrade sharply
beyond the training length (in earlier runs, perplexity rose from $4.3$ at
$1\times$ to $20.6$ at $4\times$). This is the inherited ALiBi property,
preserved when the slopes are made learnable. We regard \emph{stable
train-short / serve-long extrapolation}, rather than the modest in-distribution
gain, as the result's most consequential feature---and the bridge to the next
section.

\section{Forward-looking discussion: extrapolation and the energy footprint of
serving}
\label{sec:energy}

\paragraph{A caveat, restated.} Nothing in this section is measured. We did not
profile datacenter power, and we make no quantitative energy claim. What follows
is order-of-magnitude reasoning about \emph{where} a robust-extrapolation
positional scheme could reduce computational work, offered to motivate the
roadmap in Section~\ref{sec:roadmap}, not to assert a result.

\paragraph{Why extrapolation, specifically, touches energy.} The dominant energy
costs in the LLM lifecycle that a positional scheme can plausibly influence are:
(i) \emph{long-context adaptation}---the additional training/fine-tuning used to
extend a model's usable context (e.g.\ position-interpolation or RoPE-scaling
passes); and (ii) \emph{serving headroom}---the practice of training or
fine-tuning at a longer context than strictly necessary to guarantee quality at
the longest expected prompt. A positional scheme that degrades gracefully beyond
its training length attacks both: if a model trained at context $L$ retains
quality at $4$--$8L$ without a dedicated extension pass, the energy of that pass
is avoided, and the training context itself can be set by the \emph{typical}
rather than the \emph{worst-case} prompt.

\paragraph{An illustrative budget (not a claim).} Self-attention cost scales as
$\mathcal{O}(L^2)$ per layer in compute and---inside FlashAttention---roughly
$\mathcal{O}(L)$ in HBM traffic. A context-extension fine-tune is itself a
training job whose energy is a non-trivial fraction of a serving fleet's monthly
budget at scale. \emph{If} graceful extrapolation removed the need for such a
pass for a class of deployments, the saving would be the entire energy of that
pass; \emph{if} it allowed a fleet to train at $L$ rather than $2L$, the
quadratic term in the affected training compute would fall by up to $4\times$ on
those steps. We deliberately phrase these as conditionals: each depends on the
extrapolation holding at production scale and quality, which our $25$M-parameter
WikiText-2 evidence does \emph{not} establish.

\paragraph{What is genuinely free here.} Independently of the above, the scheme
adds essentially zero inference overhead: $H$ learnable scalars, applied as an
additive bias already fused into FlashAttention's inner loop for ALiBi. So
whatever extrapolation benefit exists comes at no measurable per-token energy
cost---unlike methods that buy context length with extra computation. This is the
honest, defensible half of the energy story: \emph{the mechanism is cheap}; the
\emph{open question} is whether its extrapolation translates into avoided
training/fine-tuning at scale.

\paragraph{Datacenter and $\mathrm{CO_2}$ framing.} Inference now dominates the
lifecycle energy of widely deployed models, and context length is a primary
driver of per-request cost. A positional scheme that is (a) free at inference and
(b) robust to serving beyond its training length is, in principle, aligned with
reducing both per-request energy variance and the amortized training energy
attributable to long-context support. Translating ``aligned in principle'' into
``kWh and $\mathrm{CO_2}$ saved'' requires the measurements we have not made. We
flag this as the single most valuable thing a reader with a serving fleet could
contribute.

\section{Lessons learned (methodological)}
\label{sec:lessons}

\begin{enumerate}[leftmargin=1.6em,itemsep=3pt]
\item \textbf{A short training budget crowns the overfitter.} An early
  $1200$-step screen ranked the most expressive bias first; at $6000$ steps the
  ranking \emph{inverted} (the expressive bias overfit a $2$\,MB corpus over
  $\sim$$37$ epochs, train loss $3\times$ lower but validation $3\times$ worse).
  Screens must validate at, or near, the target budget---or include explicit
  generalization controls.
\item \textbf{Expressive positional biases need a generalization budget, not a
  fit budget.} $\ell_2$ regularization on the learnable slopes plus validation
  early-stopping converted the inversion into a clean $+8.1\%$ PASS. Without them
  the same scheme fails.
\item \textbf{Attribution beats celebration.} Once a learnable bias beat the
  baselines, the decisive experiment was not to ship it but to ask \emph{why} it
  won. The nested control showed the win was learnability, not the geometric
  prior we started from---and saved us from publishing a false mechanism.
\item \textbf{Test the orthogonal axis, then stop.} Having exhausted ``more
  distance shape,'' we tested the one untried axis (content/position balance). It
  tied. Knowing \emph{when the search is done}---when added structure no longer
  pays---is itself a result.
\item \textbf{Negative and null results are the product.} Three of our four
  studies returned ``no improvement.'' Reported together, they are what gives the
  surviving positive result its credibility.
\end{enumerate}

\section{Vision, roadmap, and open invitations}
\label{sec:roadmap}

\paragraph{Vision.} We are interested in positional mechanisms that are
\emph{cheap, adaptable, and robust to serving beyond their training regime},
because that triple is where a small architectural choice can have an
outsized, fleet-level energy footprint. Learnable-ALiBi is a minimal instance;
the forward question is how far the ``adaptability beats added structure''
principle generalizes (to RoPE frequencies made learnable, to per-layer rather
than per-head slopes, to mixtures that remain $\mathcal{O}(1)$ state).

\paragraph{Roadmap (concrete, falsifiable).}
\begin{enumerate}[leftmargin=1.6em,itemsep=2pt]
\item \textbf{Validate on a real open-weight model.} The cheapest \emph{valid}
  test is not from-scratch but a surgical edit of an ALiBi-native open-weight
  model (frozen RoPE models cannot be tested by a positional swap---that is the
  invalid frozen-swap that produces a vacuous PASS). Concretely: take BLOOM-1b1
  (or MPT-7B at scale), make its fixed per-head slopes trainable (initialized at
  the released ladder), and run a short, parameter-efficient continued-pretrain;
  evaluate on a train-short/serve-long battery (RULER, PG-19 sliding-window
  perplexity, passkey retrieval) against the frozen fixed-ALiBi release. The
  primary metric is the \emph{usable-context multiple per unit of adaptation
  compute}, which directly tests item~2. A full specification (models, licenses,
  intervention, pre-committed kill criteria) is in
  \texttt{docs/BENCHMARK\_SPEC\_OPENWEIGHT.md}.
\item \textbf{Fair long-range bake-off (from scratch, family-neutral).}
  Learnable-ALiBi vs.\ NoPE vs.\ RoPE $+$ position-interpolation vs.\ fixed
  ALiBi at $\ge$$10\times$ the parameter scale used here, pre-committing the gate
  and the extrapolation protocol---the cleaner but costlier complement to the
  open-weight test above.
\item \textbf{Measure the energy claim or retract it.} Instrument a serving
  harness; quantify, in kWh and $\mathrm{CO_2}\text{e}$, the cost of a
  position-interpolation extension pass that graceful extrapolation could avoid.
  Report a measured number or state that no saving was found.
\item \textbf{Learnable RoPE frequencies.} Apply the same attribution discipline
  to RoPE: are learnable per-head frequency scales the RoPE analogue of our
  result, and do extra degrees of freedom again fail to survive?
\item \textbf{Exact-kernel orthogonal track.} A separate finding (rescaled-integer
  $QK^\top$ is $\sim$$34\times$ more accurate than FP16 at $L{=}2048$, bottlenecked
  by the FP16 softmax) suggests an exact-$QK^\top$ $+$ FP32-softmax kernel; pair it
  with learnable-ALiBi and re-measure quality and latency.
\end{enumerate}

\paragraph{Invitations.} We would especially value: (a) a refutation or
confirmation of the extrapolation behaviour at larger scale; (b) a fair
comparison against the baselines above run by someone without our priors; and (c)
a measured energy/$\mathrm{CO_2}$ accounting from anyone operating a serving
fleet. Code, raw run logs, and an open benchmark dataset are linked from the
repository~\cite{mms}; the full chronological arc, including the false positive we
caught, is documented there as a case study in falsification-first applied
research.

\section*{Reproducibility}
All runs use the public scripts in \texttt{4\_ai\_hardware\_attention/} of the
repository; the decisive comparison is \texttt{cy\_sieve\_hetero\_pos.py}, raw
outputs are in \texttt{gpu\_phase\_runs/hetero\_pos\_gpu\_20260623.\{json,log\}},
and the open benchmark dataset
(\href{https://huggingface.co/datasets/callensxavier/cy-sieve-attention-benchmark}{callensxavier/cy-sieve-attention-benchmark})
carries the per-scheme perplexities and the documented arc.

\begin{thebibliography}{9}
\bibitem{mms} X.~Callens, \emph{Mirror-Map-Sieve} (code, raw runs, and the
  companion preprint v7), \url{https://github.com/xaviercallens/Mirror-Map-Sieve},
  2026.
\bibitem{alibi} O.~Press, N.~A.~Smith, M.~Lewis, \emph{Train Short, Test Long:
  Attention with Linear Biases Enables Input Length Extrapolation} (ALiBi),
  ICLR 2022.
\bibitem{nope} A.~Kazemnejad et al., \emph{The Impact of Positional Encoding on
  Length Generalization in Transformers} (incl.\ NoPE), NeurIPS 2023.
\bibitem{flash} T.~Dao et al., \emph{FlashAttention: Fast and Memory-Efficient
  Exact Attention with IO-Awareness}, NeurIPS 2022.
\bibitem{rope} J.~Su et al., \emph{RoFormer: Enhanced Transformer with Rotary
  Position Embedding} (RoPE), 2021.
\bibitem{wikitext} S.~Merity, C.~Xiong, J.~Bradbury, R.~Socher, \emph{Pointer
  Sentinel Mixture Models} (WikiText), 2016.
\end{thebibliography}

\end{document}
